\paragraph{\texorpdfstring{$W_1$}{W1} 口径下的 Lee--Yang 极限测度不可统一逼近}
沿用定理 \ref{thm:fold-computability-halting-leyang-holographic-encoding} 的极限测度
\(
\mu=\mu_{\mathcal H}
\)。
在闭单位圆盘 \(\overline{\mathbb D}\subset\CC\) 上取欧氏距离诱导的一阶 Wasserstein 距离 \(W_1\)。

\begin{theorem}[Lee--Yang 极限测度在 \texorpdfstring{$W_1$}{W1} 距离下的统一逼近不可实现]\label{thm:fold-computability-halting-no-uniform-w1-learning}
\leanverified{paper\_fold\_computability\_halting\_no\_uniform\_w1\_learning}
不存在总停机算法 \(\mathcal B\)，使得对每个 \(e\ge 1\) 都输出一概率测度 \(\nu_e\in\mathrm{Prob}(\overline{\mathbb D})\) 并满足
\[
W_1(\nu_e,\mu)\le \frac{2^{-e}}{12\,q_e},
\]
其中 \(q_e\) 为定理 \ref{thm:fold-computability-halting-leyang-holographic-encoding} 所用的第 \(e\) 个奇素数。
\end{theorem}
\begin{proof}
对固定 \(e\)，取测试函数
\[
f_e(z):=z^{q_e}\qquad(z\in\overline{\mathbb D}).
\]
在 \(\overline{\mathbb D}\) 上有
\(
|f_e'(z)|\le q_e
\)，
故 \(f_e\) 是 \(q_e\)-Lipschitz。
由 Kantorovich--Rubinstein 对偶表述，
\[
\left|\int_{\overline{\mathbb D}} z^{q_e}\,d\nu_e-\int_{\overline{\mathbb D}} z^{q_e}\,d\mu\right|
\le q_e\,W_1(\nu_e,\mu)
\le \frac{2^{-e}}{12}.
\]
另一方面，由定理 \ref{thm:fold-computability-halting-leyang-holographic-encoding}，
\[
\int_{\mathbb T} z^{q_e}\,d\mu=
\frac{2^{-e}}{1+S}\,\mathbf 1_{\{e\in\mathcal H\}},
\qquad
1+S\le 2.
\]
于是
\[
e\notin\mathcal H
\Longrightarrow
\int z^{q_e}\,d\mu=0,
\qquad
e\in\mathcal H
\Longrightarrow
\int z^{q_e}\,d\mu\ge 2^{-e-1}.
\]
因此以上误差界允许用阈值 \(2^{-e}/6\) 有效区分 \(e\in\mathcal H\) 与 \(e\notin\mathcal H\)，从而判定停机集 \(\mathcal H\)，矛盾。
证毕。
\end{proof}

\begin{corollary}[极限测度在 \texorpdfstring{$W_1$}{W1} 口径下不可计算]\label{cor:fold-computability-halting-measure-not-w1-computable}
不存在总停机算法 \(\mathcal A\)，使得对每个精度参数 \(n\ge 1\) 都输出一概率测度 \(\widehat\mu_n\in\mathrm{Prob}(\overline{\mathbb D})\) 并满足
\[
W_1(\widehat\mu_n,\mu)\le 2^{-n}.
\]
亦即，\(\mu\) 在 \(W_1\) 度量下不是可计算概率测度。
\end{corollary}
\begin{proof}
若存在这样的 \(\mathcal A\)，则对每个 \(e\) 取
\[
n(e):=e+\left\lceil \log_2(12q_e)\right\rceil.
\]
于是
\[
W_1(\widehat\mu_{n(e)},\mu)\le 2^{-n(e)}\le \frac{2^{-e}}{12q_e},
\]
与定理 \ref{thm:fold-computability-halting-no-uniform-w1-learning} 矛盾。
证毕。
\end{proof}

\endinput
